
%%% Titulní strana práce

\pagestyle{empty}
\begin{center}

\large

Charles University in Prague

\medskip

Faculty of Mathematics and Physics

\vfill

{\bf\Large BACHELOR THESIS}

\vfill

\centerline{\mbox{\includegraphics[width=60mm]{img/logo.eps}}}

\vfill
\vspace{5mm}

{\LARGE Jiří Maršík}

\vspace{15mm}

% Název práce přesně podle zadání
{\LARGE\bfseries Fast and Trainable Tokenizer \\ for Natural Languages}

\vfill

% Název katedry nebo ústavu, kde byla práce oficiálně zadána
% (dle Organizační struktury MFF UK)
Institute of Formal and Applied Linguistics

\vfill

\begin{tabular}{rl}

Supervisor of the bachelor thesis: & RNDr. Ondřej Bojar, Ph.D. \\
\noalign{\vspace{2mm}}
Study programme: & Computer Science \\
\noalign{\vspace{2mm}}
Specialization: & General Computer Science \\
\end{tabular}

\vfill

% Zde doplňte rok
Prague 2011

\end{center}

\newpage

%%% Následuje vevázaný list -- kopie podepsaného "Zadání bakalářské práce".
%%% Toto zadání NENÍ součástí elektronické verze práce, nescanovat.

%%% Na tomto místě mohou být napsána případná poděkování (vedoucímu práce,
%%% konzultantovi, tomu, kdo zapůjčil software, literaturu apod.)

\openright

\noindent
Thanks for all the fish!

\newpage

%%% Strana s čestným prohlášením k bakalářské práci

\vglue 0pt plus 1fill

\noindent
I declare that I carried out this bachelor thesis independently, and only with
the cited sources, literature and other professional sources.

\medskip\noindent
I understand that my work relates to the rights and obligations under the Act
No. 121/2000 Coll., the Copyright Act, as amended, in particular the fact that
the Charles University in Prague has the right to conclude a license agreement
on the use of this work as a school work pursuant to Section 60 paragraph 1 of
the Copyright Act.

\vspace{10mm}

\hbox{\hbox to 0.5\hsize{%
In ........ date ............
\hss}\hbox to 0.5\hsize{%
signature
\hss}}

\vspace{20mm}
\newpage

%%% Povinná informační strana bakalářské práce

\vbox to 0.5\vsize{
\setlength\parindent{0mm}
\setlength\parskip{5mm}

\begin{singlespace}
\begin{otherlanguage}{czech}

Název práce:
Rychlý a trénovatelný tokenizér pro přirozené jazyky
% přesně dle zadání

Autor:
Jiří Maršík

Katedra:  % Případně Ústav:
Ústav formální a aplikované lingvistiky
% dle Organizační struktury MFF UK

Vedoucí bakalářské práce:
RNDr. Ondřej Bojar Ph.D., Ústav formální a aplikované lingvistiky
% dle Organizační struktury MFF UK, případně plný název pracoviště mimo MFF UK

Abstrakt:
% abstrakt v rozsahu 80-200 slov; nejedná se však o opis zadání bakalářské
% práce
V této práci představujeme systém pro dezambiguaci hranic mezi tokeny a větami
v textu. Charakteristickým znakem programu je jeho značná konfigurovatelnost
and všestrannost, jelikož si tokenizér dokáže poradit i s nepřerušovaným
čínským textem. Tokenizér v jádru používá klasifikátory založené na modelech s
maximální entropií a jedná se tudíž o učící se systém, kterému je nutné
předložit již tokenizovaná data k trénování. Program je doplněn nástrojem pro
hlášení úspěšnosti tokenizace, což pomáhá zejména při rychlém vývoji a ladění
tokenizačního procesu. Program používá pouze multiplatformní knihovny a při
vývoji byl kladen důraz zejména na efektivitu a správnost.
Samotná práce se skládá z přehledu dalších nástrojů na tokenizaci a segmentaci
textu a krátkého úvodu do modelů s maximální entropií. Velkou část práce pak
tvoří popis samotné implementace systému. Práce je nakonec završena
vyhodnocením úspěšnosti několika tokenizačních metod implementovaných skrz
tento systém.

Klíčová slova:
tokenizace, segmentace, maximální entropie, předzpracování textu
% 3 až 5 klíčových slov

\end{otherlanguage}
\end{singlespace}

\vss}\nobreak\vbox to 0.49\vsize{
\setlength\parindent{0mm}
\setlength\parskip{5mm}

\begin{singlespace}

Title:
Fast and Trainable Tokenizer for Natural Languages
% přesný překlad názvu práce v angličtině

Author:
Jiří Maršík

Department:
Institute of Formal and Applied Linguistics
% dle Organizační struktury MFF UK v angličtině

Supervisor:
RNDr. Ondřej Bojar Ph.D., Institute of Formal and Applied Linguistics
% dle Organizační struktury MFF UK, případně plný název pracoviště
% mimo MFF UK v angličtině

Abstract:
In this thesis, we present a data-driven system for disambiguating token and
sentence boundaries in linguistic data. The system is highly configurable and
versatile to the point its tokenization abilities are able to satisfactorily
segment unbroken Chinese text. The tokenizer relies on maximum entropy
classifiers and requires tokenized and segmented text as training data. The
program is complete with a tool for reporting the performance of the
tokenization which helps rapidly develop and tune the tokenization process. The
program is written using multiplatform libraries only and with emphasis on
speed and completeness.
The thesis itself consists of a survey of other tools for text tokenization and
segmentation and a short introduction to maximum entropy modelling. A large
part of the thesis is the description of the internal implementation. The
thesis then ends with an evaluation of several example tokenization methods
implemented using the system.
% abstrakt v rozsahu 80-200 slov v angličtině; nejedná se však o překlad zadání
% bakalářské práce

Keywords:
% 3 až 5 klíčových slov v angličtině
tokenization, segmentation, maximum entropy, text preprocessing

\end{singlespace}

\vss}

\newpage
